%/*% pdflatex -file-line-error -interaction nonstopmode %

\documentclass{article}
\usepackage{graphicx}
\usepackage{amsmath}
\usepackage{amsfonts}
\usepackage{bookmark}
\usepackage{booktabs}
\usepackage{exercises}
\graphicspath{ {./images} }
\usepackage{indentfirst}
\usepackage{amsthm}
\usepackage{yhmath}
\usepackage{placeins}
\newcommand*{\oldneg}{\mathord{\sim}}
\newcommand\showdiv[1]{\overline{\smash{)}#1}}
\newcommand\Mydiv[1]{%
\smash{\raise.3ex\hbox{$\big)$}}\mkern-8mu
        \overline{\enspace\strut#1}}
\setlength{\tabcolsep}{5pt}
\theoremstyle{definition}
\newtheorem{definition}{Definition}

\setlength{\parindent}{.5in}
\setlength{\parskip}{12pt}
\renewcommand{\arraystretch}{1}

\title{Quiz 1 (Page 135) Answers}
\begin{document}
\maketitle

\section*{1}
\[ 2x^3-54=2(x^3-27) \]
\[ =2(x-3)(x^2+3x+9) \]

\section*{2}
\[ x^3-3x^2+2x-6=x^2(x)+x^2(-3)+2(x)+2(-3) \]
\[ =(x^2+2)(x-3) \]

\section*{3}
\[ x^3+x^2+x+1=x^2(x)+x^2(1)+1(x)+1(1) \]
\[ =(x^2+1)(x+1) \]

\section*{4}
\[ 6x^5-150x=6x(x^4-25) \]
\[ =6x(x^2+5)(x^2-5) \]

\section*{5}
\[ 3x^4-24x^2+48=3(x^4-8x^2+16) \]
\[ =3(x^2-4)^2 \]

\section*{6}
\[ 2x^3-3x^2-12x+18=x^2(2x)+x^2(-3)-6(2x)-6(-3) \]
\[ =(x^2-6)(2x-3) \]

\section*{7}
\begin{center}
	\begin{tabular}{rrrrrrrrr}
		      &       &      &   &       &         & $x^2$    & $-4x$  & $+14$  \\
		                     \cmidrule{4-9}
		$x^2$ & $+5x$ & $-2$ & ) & $x^4$ & $+x^3$  & $-8x$    & $+5x$  & $+5$   \\
		      &       &      &   & $x^4$ & $+5x^3$ & $-2x^2$                   \\
			             \cmidrule{4-9}
		      &       &      &   &       & $-4x^3$ & $-6x^2$  & $+5x$  & $+5$   \\
		      &       &      &   &       & $-4x^3$ & $-20x^2$ & $+8x$          \\
						\cmidrule{6-9}
		      &       &      &   &       &         & $14x^2$  & $-3x$  & $+5$  \\
		      &       &      &   &       &         & $14x^2$  & $+70x$ & $-28$ \\
							   \cmidrule{7-9}
		      &       &      &   &       &         &          & $-73x$ & $+33$
	\end{tabular}
\end{center}
\[ \dfrac{x^4+x^3-8x^2+5x+5}{x^2+5x-2}=x^2-4x+14+\dfrac{-73x+33}{x^2+5x-2} \]

\section*{8}
\begin{center}
	\begin{tabular}{l|rrrrr}
		$-7$ & $4$ & $27$ & $3$ & $64$ \\
		& & $-28$ & $7$ & $-77$ \\
		\cmidrule{2-6}
		& $4$ & $-1$ & $11$ & $-13$
	\end{tabular}
\end{center} 
\[ \dfrac{4x^3+27x^2+3+64}{x+7}=4x^2-x+11+\frac{-13}{x+7} \]

\section*{9}
By the rational roots theorem, the possible real zeros are:
\[ \pm\frac{1}{1},\pm\frac{1}{2},\pm\frac{5}{1},\pm\frac{5}{2},\pm\frac{6}{1},\pm\frac{6}{2},\pm\frac{30}{1},\pm\frac{30}{2} \]
Using a calculator, one can see that $-\frac{1}{2}$ is one such zero. We can divide using synthetic division to find the other factor:
\begin{center}
	\begin{tabular}{l|rrrr}
		$-\frac{1}{2}$ & $2$ & $-19$ & $50$ & $30$ \\
		& & $-1$ & $10$ & $-30$ \\
		\cmidrule{2-5}
		& $2$ & $-20$ & $60$ & $0$
	\end{tabular}
\end{center}
meaning the other factor is \[ 2x^2-20x+60 \] Because \[ b^2-4ac=(-20)^2-4(2)(60)=-80<0 \], the other zeros of the function are not real, meaning 
$-\frac{1}{2}$ is the only zero. 

\section*{10}
By the rational roots theorem, the possible zeros are:
\[ \pm\frac{1}{1},\pm\frac{2}{1},\pm\frac{4}{1},\pm\frac{7}{1},\pm\frac{8}{1},\pm\frac{14}{1},\pm\frac{28}{1},\pm\frac{56}{1} \]
Using a calculator, one can see that $8$ is one such zero. We can divide using synthetic division to find the other factor:
\begin{center}
	\begin{tabular}{l|rrrr}
		$8$ & $1$ & $-4$ & $-25$ & $-56$ \\
		& & $8$ & $32$ & $56$ \\
		\cmidrule{2-5}
		& $1$ & $4$ & $7$ & $0$
	\end{tabular}
\end{center}
meaning the other factor is \[ x^2+4x+7 \] Because $b^2-4ac=4^2-4(1)(7)=-12<0 \] there are no other real zeros. to the function, meaning 
$8$ is the only real zero of the function. 

\section*{13}


We are designing a square patio that is to be made from 128 cubic feet. The thickness of the patio is 15.5 ft. less than each side length, meaning 
we have the following equation where $x$ is the side length:
\[ 128=x^2(x-15.5) \]
\[ x^3-15.5x^2-128=0 \]

\end{document}
